% Abstract v3 - convergence thesis as the centerpiece (hat-trick protagonist).
\begin{abstract}
We ask how much of engineered network load balancing is recoverable from
classical physics alone, and find that the answer does not depend on which
physics one picks. We formalise three independent routing cores, each a packet
treated as a particle under a different classical law: a Newtonian
reaction (a dropped packet leaves a decaying scar that repels its successors),
an Archimedean buoyancy (edges push packets away non-linearly once congestion
crosses a flotation threshold), and a Pascalian pressure that diffuses
congestion to neighbouring edges. The three share a single search engine and no
mechanism. Tested as statistical equivalence (TOST, margin $2$ percentage
points) across fifteen topologies against the engineered balancers CONGA (at a
fair path budget) and DRILL, all three reach parity in the same band---nine to
twelve of fifteen topologies---and, more tellingly, they succeed and fail on the
\emph{same} topologies. A purely structural metric, the tube/sp ratio (room to
manoeuvre among near-shortest paths, computable before deployment), predicts the
shared failures of all three cores: the two real WAN backbones, with low
tube/sp, are exactly where engineered multi-path routing keeps an edge. That
independent classical-physics mechanisms converge on the same operating point,
with the same boundary, indicates that the result is a property of the load-
balancing problem rather than of any one analogy: where a topology offers room
to manoeuvre, classical physics recovers engineered performance with a fraction
of the machinery, and a single ratio says where. We report negative results, a
self-audit of an earlier inflated finding, and the per-topology disagreements in
full.
\end{abstract}
